% SOURCE AUTHORITY:
% expose_VII_body.tex lines 1545--1662; scan indices 209--215; source folios 198--204.
% Direct scan comparison: physical PDF pages 210--216 at 600 dpi.
% Transparent corrections: the first overbrace in (3.5.1) is L_2 as printed, not
% the working transcription's L_o; chain subscripts o are rendered as zero; the
% chain-complex relation is written in compositional order d_0d_1=0 throughout;
% and the display used later as (*) receives the tag visible in the scan.
% Hard stop before source line 1668, section 3.6.

\subsection*{3.5. A canonical partial flat resolution of an
arbitrary $\mathbb Z$-module}

We apply the preceding results when $\underline C$ is the category
of commutative groups of the topos $\underline T$, that is, the
category of $\mathbb Z$-modules of $\underline T$.  Let $P$ be an
object of $\underline C$.  We shall define, functorially in $P$, a
chain complex $L_\bullet(P)$ augmented to $P$:
\begin{equation}\tag{3.5.1}
\begin{tikzcd}[column sep=small,row sep=large]
 0\arrow[r]&\overbrace{\mathbb Z[P\times P]\times
   \mathbb Z[P\times P\times P]}^{L_2}
   \arrow[r,"d_1"]
 &\overbrace{\mathbb Z[P\times P]}^{L_1}
   \arrow[r,"d_0"]
 &\overbrace{\mathbb Z[P]}^{L_0}
   \arrow[r]\arrow[d,"\varepsilon"']&0\\
 &&&P&{}
\end{tikzcd}
\end{equation}
with $L_i(P)=0$ for $i\geq3$.  The components are displayed in
(3.5.1), using the notation of 1.4.  To describe the differentials
$d_0,d_1$ and the augmentation $\varepsilon$, write $[p]$ for the
section of $\mathbb Z[P](S)$ defined by $p\in P(S)$, and similarly
write $[p,q]$ and $[p,q,r]$ for the corresponding sections of
$\mathbb Z[P\times P](S)$ and
$\mathbb Z[P\times P\times P](S)$.  Then set
\begin{equation}\tag{3.5.2}
\left\{
\begin{aligned}
 \varepsilon[p]&=p,\\
 d_0[p,q]&=[p+q]-[p]-[q],\\
 d_1[p,q]&=[p,q]-[q,p],\\
 d_1[p,q,r]&=[p+q,r]-[p,q+r]+[p,q]-[q,r].
\end{aligned}
\right.
\end{equation}
These homomorphisms of commutative groups do define a complex
$L_\bullet(P)$ augmented to $P$:
\[
 \varepsilon d_0=0,\qquad d_0d_1=0.
\]
The first relation is immediate.  The second splits into two
relations according to the two factors of $L_2(P)$; they express,
respectively, commutativity and associativity of the composition law
on $P$.

The augmented complex depends functorially on $P$.  A morphism
$u:P\longrightarrow P'$ gives a commutative diagram
\[
\begin{tikzcd}[column sep=huge,row sep=large]
 L_\bullet(P)\arrow[r,"L_\bullet(u)"]\arrow[d]
 &L_\bullet(P')\arrow[d]\\
 P\arrow[r,"u"]&P'.
\end{tikzcd}
\]
Its components are flat, as is every free $\mathbb Z$-module.

\underline{Proposition} 3.5.3. \underline{One has}
$H_1(L_\bullet(P))=0$, \underline{and the augmentation}
$\varepsilon$ \underline{induces an isomorphism}
\[
 H_0(L_\bullet(P))\xrightarrow{\sim}P.
\]

For the proof one can make the standard reduction to the case in
which $\underline T$ is the category of sets and $P$ is an ordinary
$\mathbb Z$-module, where (3.5.1) is likely familiar.  We instead
give a direct proof that does not use this reduction, applying 3.4.7
to the augmentation
\[
 \varepsilon:L_\bullet(P)\longrightarrow P,
\]
where $P$ is regarded, as usual, as a chain complex concentrated in
degree zero.  It remains to prove that, for every object $G$ of
$\underline C$, the functor
\begin{equation}\tag{*}
 \varepsilon^*:\underline\Psi_P(G)\longrightarrow
 \underline\Psi_{L_\bullet(P)}(G)
\end{equation}
is an equivalence of categories.

The source of (*) is plainly isomorphic to the category of
extensions of the $\mathbb Z$-module $P$ by $G$.  To describe its
target, use 1.4 to identify the categories of extensions of
$L_0,L_1,L_2$ by $G$, respectively, with the category of torsors
under $G_P$, the category of torsors under $G_{P\times P}$, and the
category of pairs consisting of a torsor under $G_{P\times P}$ and
a torsor under $G_{P\times P\times P}$.

Thus an object of the target of (*) is a pair $(E,\varphi)$, where
$E$ is a torsor under $G_P$ and $\varphi$ is a trivialization of the
torsor $E'$ on $P\times P$ corresponding to the inverse image along
$d_0$ of the extension of $\mathbb Z[P]$ by $G$ defined by $E$.
The trivialization is subject to the compatibility with
$d_0d_1=0$ described below.  Formula (3.5.2) immediately gives
\[
 E'=\pi^*E\,\pr_1^*(E)^{-1}\,\pr_2^*(E)^{-1},
\]
and a trivialization of this torsor is the same as an isomorphism,
also denoted by $\varphi$,
\[
 \varphi:\pr_1^*E\,\pr_2^*E\xrightarrow{\sim}\pi^*E.
\]
Expanding the compatibility of $\varphi$ with $d_0d_1=0$ gives the
commutativity condition (1.2.1) and the associativity condition
(1.1.4.1), corresponding to the two factors of $L_2$.  The assertion
that (*) is an equivalence is therefore precisely the dictionary of
1.1.6.  This proves 3.5.3.

\underline{Remark} 3.5.4. a) It would be interesting to strengthen
3.5.3 by constructing a resolution $L_\bullet(P)$ of $P$, functorial
in $P$, whose components are functorially isomorphic to sums of free
$\mathbb Z$-modules generated by cartesian powers of $P$.  In view
of isomorphisms of the form (1.4.5), this would give a method for
partially computing $\Ext^i(P,A)$ as the abutment of a spectral
sequence whose initial term is expressed through groups
\[
 H^j(P\times P\times\cdots\times P,A).
\]
Here we obtain only very partial information about $\Ext^1$, and no
information about the higher $\Ext^i$.  There is, in any case, a
natural way to continue the partial resolution $L_\bullet(P)$ by one
degree.\footnote{L. Breen constructed such a resolution (editor's
note).}

b) There is a variant of (3.5.1) and 3.5.3 in the category of
$A$-modules, where $A$ is a fixed ring of $\underline T$.  The form
of $L_\bullet(P)$ is more complicated because the description 1.1.6
must be enlarged (compare 2.10.3).  One obtains
\begin{equation}\tag{3.5.5}
\left\{
\begin{aligned}
 L_0={}&A[P],\\
 L_1={}&A[P\times P]+A[A\times P],\\
 L_2={}&A[P\times P\times P]+A[P\times P]
       +A[A\times A\times P]+A[A\times P\times P].
\end{aligned}
\right.
\end{equation}
The augmented-complex structure is given by (3.5.2), supplemented by
\begin{equation}\tag{3.5.6}
\left\{
\begin{aligned}
 d_0[\lambda,p]={}&[\lambda p]-\lambda[p],\\
 d_1[\lambda,\mu,p]={}&[\lambda,\mu p]-[\lambda\mu,p]
                         +\lambda[\mu,p],\\
 d_1[\lambda,p,q]={}&[\lambda,p+q]-[\lambda,p]-[\lambda,q]
                       +\lambda[p,q]-[\lambda p,\lambda q].
\end{aligned}
\right.
\end{equation}
The relation $d_0d_1=0$ follows from the analogous relation on the
components appearing in (3.5.2), which, as we saw, is equivalent to
associativity and commutativity of addition on $P$, together with
the relations
\[
 d_0d_1[\lambda,\mu,p]=0,
 \qquad d_0d_1[\lambda,p,q]=0.
\]
The latter express, respectively, associativity of the given action
of $A$ on $P$,
\[
 \lambda(\mu p)=(\lambda\mu)p,
\]
and its distributivity over addition,
\[
 \lambda(p+q)=\lambda p+\lambda q.
\]

For a torsor $E$ under $G_P$, a trivialization of $d_0^*E$ amounts
to an internal composition law on $E$ (satisfying the compatibility
conditions already described for the internal law on $E$ and its
torsor structure), together with an external composition law of $A$
on $E$, compatible with the action of $A$ on $P$ and on $G$.  The
compatibility of this trivialization with the evident trivialization
of $d_1^*d_0^*E$ is equivalent to the conditions already described
above, which say that the internal law on $E$ is a commutative-group
law extending $P$ by $G$, together with associativity of the external
law and distributivity of the latter with respect to the internal
law.  These conditions say exactly that $E$, with its two
composition laws and its torsor structure, is an $A$-module extension
of the $A$-modules $P$ and $G$.

It follows, by essentially the same method as in 3.5.3, that
$L_\bullet$ is a partial resolution of the module $P$, that is,
$H_1(L_\bullet)=0$.  It would again be interesting to extend this
partial resolution to an infinite resolution of the same kind, with
components of the form
\[
 A[A\times A\times\cdots\times P\times P\times\cdots].
\]
Compare the analogous question in a) for commutative groups.
